\clearpage
\begin{explanationbox}
\textbf{第一层：先看两段长度——为什么是 $a\pm A$？}

\textbf{先看具体数值：}若两条焦半径总长是 $14$，长短之差是 $4$，
则平均长度是 $7$，较长的比平均值多 $2$，较短的比平均值少 $2$。

\begin{center}
% C093：长度分解教学图；真实线段长度与 9:7:5:2 一致。
% 长焦半径=9=7+2，短焦半径=5=7-2。
\begin{tikzpicture}[x=0.76cm,y=0.88cm,>=stealth]
  % 平均值分界 7
  \draw[gray!65,densely dashed] (7,-0.9)--(7,2.65);
  \node[gray!70!black,above] at (7,2.65) {$a=7$};
  \node[anchor=west] at (-0.4,3.2) {\textbf{同一平均值 $7$，一长一短}};
  % Long bar 0..7 plus 7..9
  \draw[teal!70!black,line width=5pt] (0,1.9)--(7,1.9);
  \draw[orange!88!black,line width=5pt] (7,1.9)--(9,1.9);
  \draw[black,thin] (0,1.62)--(0,2.18) (9,1.62)--(9,2.18);
  \node[anchor=east] at (-0.3,1.9) {较长};
  \node[above,teal!70!black] at (3.4,1.97) {$a=7$};
  \node[above,orange!88!black] at (8.0,1.97) {$A=2$};
  \node[anchor=west] at (9.3,1.9) {$9=7+2$};
  % Short bar 0..5, dashed deficit to avg length=7
  \draw[teal!70!black,line width=5pt] (0,0.15)--(5,0.15);
  \draw[orange!85!black,densely dashed,line width=1.7pt] (5,0.15)--(7,0.15);
  \draw[black,thin] (0,-0.12)--(0,0.43) (5,-0.12)--(5,0.43);
  \node[anchor=east] at (-0.3,0.15) {较短};
  \node[above,teal!70!black] at (2.4,0.22) {$a-A=5$};
  \node[below,orange!88!black] at (6,-0.05) {少 $2$};
  \node[anchor=west] at (9.3,0.15) {$5=7-2$};
  \node[gray!75!black,anchor=west] at (0,-1.5) {实线是真实长度，虚线表示“距平均值还差 $2$”。};
\end{tikzpicture}

\end{center}

\textbf{再从图推广到字母：}设 $r_1=PF_1$，$r_2=PF_2$。
椭圆定义告诉我们\emph{两条焦半径加起来有多长}：
\[r_1+r_2=2a.\]
双曲线定义告诉我们\emph{两条焦半径相差多长}：
\[|r_1-r_2|=2A.\]
所以两条长度的平均值是 $a$，长短之间差 $2A$；
较长的一条高于平均值 $A$，较短的一条低于平均值 $A$。
\[
\boxed{\text{长焦半径}=a+A},\qquad
\boxed{\text{短焦半径}=a-A}.
\]
这里的“相加、相减”是从\emph{总长与长度差还原每一条长度}，
不是要求先解两条曲线的交点坐标。
\end{explanationbox}

\clearpage
\begin{explanationbox}
\textbf{第二层：左右支交换的是焦点的对应关系，长度本身不变。}

\begin{center}
% Two exact (3,4,5) triangles with c=2, a=4, A=1; left/right swap.
\begin{tikzpicture}[scale=0.69,>=stealth]
 \begin{scope}
  \coordinate (L) at (0,0); \coordinate (R) at (4,0); \coordinate (P) at (4,3);
  \draw[gray,thick] (L)--(R);
  \draw[orange!85!black,very thick] (L)--(P);
  \draw[teal!70!black,very thick] (R)--(P);
  \fill (L) circle (2pt) node[below]{$F_1$};
  \fill (R) circle (2pt) node[below]{$F_2$};
  \fill (P) circle (2.3pt) node[above]{$P_R$};
  \node[orange!85!black] at (1.4,2.18) {$5$};
  \node[teal!70!black] at (4.32,1.5) {$3$};
  \node[align=center] at (2,-1.1) {$x>0$\\右支：$PF_1>PF_2$};
 \end{scope}
 \begin{scope}[xshift=7.0cm]
  \coordinate (L) at (0,0); \coordinate (R) at (4,0); \coordinate (P) at (0,3);
  \draw[gray,thick] (L)--(R);
  \draw[teal!70!black,very thick] (L)--(P);
  \draw[orange!85!black,very thick] (R)--(P);
  \fill (L) circle (2pt) node[below]{$F_1$};
  \fill (R) circle (2pt) node[below]{$F_2$};
  \fill (P) circle (2.3pt) node[above]{$P_L$};
  \node[teal!70!black] at (-0.35,1.5) {$3$};
  \node[orange!85!black] at (2.55,2.18) {$5$};
  \node[align=center] at (2,-1.1) {$x<0$\\左支：$PF_1<PF_2$};
 \end{scope}
\end{tikzpicture}

\end{center}

在右支，点 $P$ 靠近右焦点 $F_2$，所以
$PF_1=a+A$、$PF_2=a-A$；
在左支，点 $P$ 靠近左焦点 $F_1$，两个公式交换。

\textbf{为什么远近只由 $x_P$ 的正负决定？}
因为 $F_1=(-c,0)$、$F_2=(c,0)$，距离公式给出
\[
\begin{aligned}
PF_1^2-PF_2^2
&=(x_P+c)^2+y_P^2-\bigl[(x_P-c)^2+y_P^2\bigr]\\
&=\boxed{4cx_P}.
\end{aligned}
\]
两边都出现的 $y_P^2$ 被减掉，结果只剩 $x_P$。
$c>0$，所以 $x_P>0$ 时 $PF_1>PF_2$；
$x_P<0$ 时 $PF_1<PF_2$。
改变 $y_P$ 的正负不会改变任一焦半径。

\textbf{判断顺序：}先算出 $a+A$、$a-A$；
再判断公共点位于左支还是右支，最后给 $PF_1,PF_2$ 分配长度。
\end{explanationbox}

\clearpage
\begin{explanationbox}
\textbf{第三层：跟离心率和其他二级结论衔接。}

\textbf{题目若只给离心率：}从
\[
e_1=\frac c{a},\quad e_2=\frac c{A}
\quad\Longrightarrow\quad a=\frac c{e_1},\quad A=\frac c{e_2}
\]
先求两种半轴，再按“平均值加减偏离量”即可。

\textbf{为什么长度都大于零？}由 $c^2=a^2-b^2=A^2+B^2$ 知 $a>c>A>0$，
必有 $a-A>0$，不会得到负长度。

\textbf{联系初中的三角形不等式：}
\begin{center}
% Exact focal triangle: a=4, c=2, A=1; side lengths 5,3,4.
\begin{tikzpicture}[x=1cm,y=1cm,font=\small]
 \coordinate (L) at (0,0); \coordinate (R) at (4,0); \coordinate (P) at (4,3);
 \draw[orange!85!black,very thick] (L)--(P)
   node[midway,above left]{$a+A=5$};
 \draw[teal!70!black,very thick] (R)--(P)
   node[midway,right=4pt]{$a-A=3$};
 \draw[gray!80,thick] (L)--(R) node[midway,below=4pt]{$2c=4$};
 \fill (L) circle (2pt) node[below left]{$F_1$};
 \fill (R) circle (2pt) node[below right]{$F_2$};
 \fill (P) circle (2pt) node[above right]{$P$};
 \node at (2,4.0) {\textbf{焦点三角形：两边之差 $<$ 底边 $<$ 两边之和}};
 \node at (2,-1.25) {$2A=2\quad <\quad 2c=4\quad <\quad 2a=8$};
\end{tikzpicture}

\end{center}
由图中的 $2<4<8$ 推广：两条焦半径之差 $2A$ 小于焦距 $2c$，
它们的和 $2a$ 大于焦距，即 $a>c>A$。

\textbf{内容分工：}C093 的重点是\emph{不先求坐标就求出焦半径}；
C094 专门研究\emph{四个公共点的坐标}；
C090 在已知焦半径的基础上，继续用余弦定理研究焦角。
\end{explanationbox}
